% ============================================================
\chapter{Hardware and Inference-Speed Analysis}
\label{ch:hardware}
% ============================================================

\section{Why One FPS Number Is Insufficient}

Inference speed is a property of a complete execution configuration, not only
of a neural-network architecture. Processor, runtime, numerical precision,
batch size, tensor shape, compiler support, memory transfer, warm-up, and the
timer boundaries all affect the result. This chapter therefore keeps three
questions separate:
\begin{enumerate}
  \item how fast does the saved model execute on a consumer CPU?
  \item how do all nine architectures compare on one common server GPU?
  \item what throughput is obtained after compiling the two deployed
  checkpoints into fixed TensorRT engines and including preprocessing and mask
  return?
\end{enumerate}

No speed value is borrowed from an architecture paper for ranking. Published
Cityscapes or A6000 FPS figures use a different image size, model output, runtime, and
hardware. Every table below names its own scope.

\section{Ryzen 5 5500 CPU Results}

Table~\ref{tab:cpu_comparison} reports batch-one forward throughput with six
PyTorch threads. Eager PIDNet-S reaches 21.16 FPS with 47.25 ms mean latency
for the reference checkpoint. Its median is 46.45 ms and P95 is 55.63 ms.
This confirms GPU-free operation at a modest perception rate, but the timer
does not include camera decode or planning.

\begin{table}[H]
\centering
\caption{Ryzen~5~5500 batch-one throughput at $640\times384$. Values are
forward and prediction only; compiled timing excludes compilation warm-up.}
\label{tab:cpu_comparison}
\small
\begin{tabular}{@{}lrrr@{}}
\toprule
\textbf{Model} & \textbf{Eager FPS} & \textbf{Compiled FPS} & \textbf{Compiled mean latency} \\
\midrule
DDRNet-23-Slim & 27.74 & \textbf{36.50} & 27.40 ms \\
PIDNet-S & 21.16 & 27.04 & 36.98 ms \\
BiSeNetV2 & 11.48 & 16.34 & 61.20 ms \\
FPN/MobileNetV2 & 7.23 & 14.26 & 70.13 ms \\
FPN/EfficientNet-B0 & 5.92 & 11.70 & 85.47 ms \\
U-Net/MobileNetV2 & 7.10 & 10.15 & 98.52 ms \\
U-Net/EfficientNet-B0 & 6.96 & 9.66 & 103.52 ms \\
SegFormer-B0 & 4.38 & 8.23 & 121.51 ms \\
ROD ViT-S & 0.41 & 0.46 & 2,195.96 ms \\
\bottomrule
\end{tabular}
\end{table}

Compiled DDRNet-23-Slim is the CPU throughput leader at 36.50 FPS. Its
convergence-aware blue+green mIoU, however, is the lowest in the nine-model
table at 0.9095. PIDNet-S is slower but reaches 0.9195.
FPN/EfficientNet-B0 reaches only 11.70 compiled FPS but gives the highest mean
accuracy, 0.9423. The CPU therefore presents a clear sequence of trade-offs
rather than one winner. ROD reaches 0.41 eager FPS and 0.46 compiled FPS; its
frozen encoder reduces training-time optimisation but the full internally
resized $1024\times1024$ transformer remains active during inference.

Compiler benefit is architecture-dependent. FPN/MobileNetV2 nearly doubles its
throughput, and SegFormer-B0 improves substantially, while U-Net and PIDNet
receive smaller gains. These differences show why floating-point-operation
(FLOP) counts or parameter estimates cannot substitute for runtime
measurements. Compilation raises ROD throughput by only 11.2\%, from 0.4094 to
0.4554 FPS, leaving the current implementation unsuitable for CPU-rate
perception loops.

\section{Common H100 Results}

The H100 protocol controls hardware and eager FP32 execution across all nine
systems. PIDNet-S is fastest at 270.38 FPS. DDRNet-23-Slim follows at 245.17
FPS, while ROD is slowest at 21.79 FPS. The frozen encoder reduces ROD training
gradients but does not reduce inference work; the full transformer remains on
the forward path.

\begin{table}[H]
\centering
\caption{Parameter count and common H100 NVL 3g.47GB MIG batch-one
forward throughput. External input is $640\times384$; execution is eager FP32.}
\label{tab:h100_comparison}
\small
\begin{tabular}{@{}lrrr@{}}
\toprule
\textbf{Model} & \textbf{Parameters} & \textbf{Mean FPS} & \textbf{Mean latency} \\
\midrule
PIDNet-S & 7.62 M & \textbf{270.38} & 3.70 ms \\
DDRNet-23-Slim & 5.69 M & 245.17 & 4.08 ms \\
BiSeNetV2 & 3.34 M & 167.80 & 5.96 ms \\
FPN/MobileNetV2 & 4.22 M & 147.36 & 6.79 ms \\
SegFormer-B0 & 3.71 M & 141.23 & 7.08 ms \\
U-Net/MobileNetV2 & 6.63 M & 131.86 & 7.58 ms \\
FPN/EfficientNet-B0 & 5.76 M & 131.40 & 7.61 ms \\
U-Net/EfficientNet-B0 & 6.25 M & 101.24 & 9.88 ms \\
ROD ViT-S & 29.11 M & 21.79 & 45.89 ms \\
\bottomrule
\end{tabular}
\end{table}

Table~\ref{tab:h100_comparison} reports the measured H100 throughput and
latency for every architecture under this common scope.

On the same H100 protocol, PIDNet-S is approximately 12.4 times faster than
ROD. Under convergence-aware training, ROD's blue+green mean is 1.36 points
higher, so that accuracy improvement has a large latency price. In the
nine-model table, ROD is also slower
and less accurate than FPN/EfficientNet-B0 on this particular hardware and
recipe. This does not make ROD universally inferior: a model-specific compiler,
different training objective, memory constraint, or another dataset can change
the decision. It does mean that the controlled H100 evidence does not place
ROD on the observed accuracy--speed frontier.

Parameter count again fails to determine speed. SegFormer-B0 has only 3.71
million parameters but is much slower than PIDNet-S on CPU. BiSeNetV2 is the
smallest system but not the fastest H100 model. Operator mix, feature-map size,
attention, memory traffic, and kernel implementation matter alongside stored
weights.

\section{Accuracy--Speed Relationship}

Figure~\ref{fig:accuracy_speed_h100} combines convergence-aware blue+green mean
mIoU with the existing H100 throughput measurements. Moving upward improves
accuracy; moving right improves speed. No model occupies the upper-right
corner.

\begin{figure}[H]
\centering
\begin{tikzpicture}
\begin{axis}[
  width=0.92\textwidth,
  height=8.0cm,
  xmode=log,
  xmin=18,xmax=320,
  ymin=0.905,ymax=0.946,
  xlabel={H100 eager FP32 throughput (FPS, log scale)},
  ylabel={Blue+green mean mIoU},
  grid=both,
  grid style={line},
  xtick={20,30,50,100,150,200,300},
  xticklabels={20,30,50,100,150,200,300}
]
\addplot[only marks,mark=*,mark size=3pt,color=azulUC3M] coordinates {
  (270.38,0.9195)
  (245.17,0.9095)
  (167.80,0.9149)
  (147.36,0.9372)
  (141.23,0.9341)
  (131.86,0.9332)
  (131.40,0.9423)
  (101.24,0.9409)
  (21.79,0.9331)
};
\node[font=\scriptsize,anchor=south west] at (axis cs:270.38,0.9195) {PID};
\node[font=\scriptsize,anchor=north east] at (axis cs:245.17,0.9095) {DDR};
\node[font=\scriptsize,anchor=north] at (axis cs:167.80,0.9149) {BiSe};
\node[font=\scriptsize,anchor=south east] at (axis cs:147.36,0.9372) {FPN-M};
\node[font=\scriptsize,anchor=north] at (axis cs:141.23,0.9341) {Seg};
\node[font=\scriptsize,anchor=north east] at (axis cs:131.86,0.9332) {UNet-M};
\node[font=\scriptsize,anchor=south west] at (axis cs:131.40,0.9423) {FPN-E};
\node[font=\scriptsize,anchor=north east] at (axis cs:101.24,0.9409) {UNet-E};
\node[font=\scriptsize,anchor=south west] at (axis cs:21.79,0.9331) {ROD};
\end{axis}
\end{tikzpicture}
\caption[Blue+green accuracy versus H100 throughput.]{Blue+green mean mIoU
versus common H100 throughput. E and M denote EfficientNet-B0 and
MobileNetV2. FPN/EfficientNet-B0 supplies the highest accuracy, while PIDNet-S
provides the highest throughput. Accuracy uses the convergence-aware
blue+green suite; inference speed uses the unchanged architecture graphs.}
\label{fig:accuracy_speed_h100}
\end{figure}

For a latency-first H100 deployment, PIDNet-S offers the most headroom. The FPN
with MobileNetV2 adds 1.77 mIoU points while retaining 147.36 FPS. Replacing
that encoder with EfficientNet-B0 adds another 0.50 points at 131.40 FPS and
gives the highest observed accuracy mean. The EfficientNet-B0 U-Net is 0.13
points lower and also has lower throughput in this runtime. The plotted means
do not include latency variance or model-specific compiler tuning, so the
frontier is specific to the declared protocol.

The CPU frontier differs. Compiled DDRNet supplies the only measured rate above
30 FPS, PIDNet-S offers a stronger accuracy point at 27.04 FPS, and
FPN/EfficientNet-B0 provides the highest accuracy at 11.70 FPS. A robot with a
10 Hz perception requirement may select a different model from one requiring
30 Hz. Frame rate should therefore be treated as a constraint in the model
choice, not as an isolated leaderboard.

\section{RTX 5060 TensorRT Results}

TensorRT changes the deployment picture for the two exported checkpoints.
Table~\ref{tab:tensorrt_speed} reports both device and measured-pipeline
scopes. PIDNet-S reaches 106.43 pipeline FPS with 9.396 ms mean latency. ROD
reaches 39.38 pipeline FPS with 25.392 ms mean latency. ROD therefore crosses
a nominal 30 FPS threshold for the measured pipeline after compilation,
although the unmeasured camera and planning stages still consume part of a
real 33.3 ms frame budget.

\begin{table}[H]
\centering
\caption[TensorRT batch-one timing on the RTX 5060.]{FP16-enabled
mixed-precision TensorRT batch-one timing on the
RTX~5060. Each distribution contains 320 measured samples after warm-up.}
\label{tab:tensorrt_speed}
\small
\begin{tabular}{@{}llrrrr@{}}
\toprule
\textbf{Model} & \textbf{Scope} & \textbf{Mean} & \textbf{P50} & \textbf{P95} & \textbf{FPS} \\
\midrule
\multirow{2}{*}{PIDNet-S}
 & Device & 0.963 ms & 0.880 ms & 1.414 ms & 1,037.97 \\
 & Measured pipeline & 9.396 ms & 9.178 ms & 10.943 ms & \textbf{106.43} \\
\midrule
\multirow{2}{*}{ROD ViT-S}
 & Device & 17.401 ms & 17.397 ms & 17.669 ms & 57.47 \\
 & Measured pipeline & 25.392 ms & 25.255 ms & 26.580 ms & \textbf{39.38} \\
\bottomrule
\end{tabular}
\end{table}

Figure~\ref{fig:tensorrt_frame_budget} expresses the same result as a frame
budget. At 30 FPS, one complete cycle may occupy at most 33.33 ms. The coloured
segment is the measured TensorRT pipeline; the pale segment is only remaining
time, not a measurement of the omitted camera, transport, planning, and control
stages.

\begin{figure}[H]
\centering
\resizebox{0.92\textwidth}{!}{\begin{tikzpicture}[font=\sffamily\small]
  \def\scale{0.39}
  \def\budget{13.0}
  \def\pid{3.664}
  \def\rod{9.903}

  \node[anchor=east,font=\bfseries] at (-0.25,2.35) {PIDNet-S};
  \node[anchor=east,font=\bfseries] at (-0.25,0.85) {ROD ViT-S};

  \fill[pidP!72] (0,1.85) rectangle (\pid,2.85);
  \fill[good!26] (\pid,1.85) rectangle (\budget,2.85);
  \draw[very thick] (0,1.85) rectangle (\budget,2.85);
  \node[text=white,font=\bfseries\footnotesize,align=center] at ({\pid/2},2.35)
    {measured pipeline\\9.396 ms};
  \node[font=\footnotesize,align=center] at ({(\pid+\budget)/2},2.35)
    {remaining 30-FPS budget\\23.937 ms};

  \fill[pidD!72] (0,0.35) rectangle (\rod,1.35);
  \fill[good!26] (\rod,0.35) rectangle (\budget,1.35);
  \draw[very thick] (0,0.35) rectangle (\budget,1.35);
  \node[text=white,font=\bfseries\footnotesize,align=center] at ({\rod/2},0.85)
    {measured pipeline 25.392 ms};
  \node[font=\footnotesize,align=center] at ({(\rod+\budget)/2},0.85)
    {remaining\\7.941 ms};

  \draw[safety,very thick,dashed] (\budget,-0.05) -- (\budget,3.25);
  \node[anchor=west,font=\footnotesize,text=safety] at (\budget+0.12,3.05)
    {33.33 ms = 30 FPS};
  \node[anchor=west,font=\footnotesize] at (0,-0.30) {0 ms};
  \node[anchor=east,font=\footnotesize] at (\budget,-0.30) {33.33 ms};
\end{tikzpicture}
}
\caption[TensorRT latency within a 30-FPS frame budget.]{Measured RTX~5060
TensorRT pipeline latency within a nominal 30-FPS frame budget. PIDNet-S leaves
substantially more unallocated time than ROD, but neither bar measures the
remaining robot stack. Original graphic drawn from the results in
Table~\ref{tab:tensorrt_speed}.}
\label{fig:tensorrt_frame_budget}
\end{figure}

The gap between PIDNet's 0.963 ms device time and 9.396 ms pipeline time shows
that host preprocessing and transfers dominate after the small network is
compiled. ROD still spends most of its time in the engine, because the internal
$1024\times1024$ transformer remains expensive. Improving CPU resize or transfer would
therefore benefit PIDNet proportionally more than ROD.

The TensorRT values should not be interpreted as a controlled speed-up over
the H100 or Ryzen tables because the GPU, runtime, and precision all change.
They answer a deployment question: what rate do the saved checkpoints achieve
after hardware-specific compilation on the available local GPU?

\section{Mixed-Precision Accuracy Verification}

Full-test evaluation confirms that mixed-precision conversion does not materially change
the operating points. Table~\ref{tab:tensorrt_accuracy} compares each engine
with its source PyTorch checkpoint.

\begin{table}[H]
\centering
\caption{Full 544-image comparison of each source PyTorch checkpoint with its
mixed-precision TensorRT conversion.}
\label{tab:tensorrt_accuracy}
\small
\begin{tabular}{@{}lrrrr@{}}
\toprule
\textbf{Model/runtime} & \textbf{mIoU} & \textbf{FSR} & \textbf{FBR} & \textbf{mIoU change} \\
\midrule
PIDNet-S PyTorch & 0.893941 & 4.3176\% & 7.0369\% & reference \\
PIDNet-S TensorRT & 0.893764 & 4.3117\% & 7.0667\% & $-0.000177$ \\
\midrule
ROD PyTorch paper recipe & 0.923569 & 3.85\% & 3.97\% & reference \\
ROD TensorRT & 0.923635 & 3.8508\% & 3.9574\% & $+0.000066$ \\
\bottomrule
\end{tabular}
\end{table}

PIDNet changes by 0.0177 mIoU percentage points and ROD by 0.0066 points.
Small positive or negative differences are expected when probabilities close
to 0.50 cross the threshold after mixed-precision execution. The full confusion
matrices, rather than a small image sample, support the claim that both
conversions preserve accuracy at the precision used for architectural
interpretation.

\section{Deployment Meaning}

Three distinct deployment candidates emerge:
\begin{itemize}
  \item \textbf{CPU-constrained:} compiled DDRNet maximises frame rate, while
  PIDNet-S offers a more accurate near-real-time alternative. The
  EfficientNet systems require a lower frame-rate budget, and the current ROD
  implementation is not a practical CPU candidate.
  \item \textbf{Powerful common GPU:} PIDNet-S maximises throughput;
  FPN/EfficientNet-B0 provides the highest accuracy above 100 FPS on the H100
  protocol, followed by U-Net/EfficientNet-B0.
  \item \textbf{Available RTX TensorRT path:} the reference PIDNet model
  exceeds 100 pipeline FPS, while the more accurate reconstructed ROD model
  reaches 39.4 FPS.
\end{itemize}

None of these measurements includes sensor exposure, decode, inter-process
transport, projection into vehicle coordinates, temporal filtering, path
planning, or control. Power draw, thermal throttling, and embedded-GPU
behaviour are also unmeasured. The results establish inference capacity and an
accuracy--latency trade-off, not production readiness of a complete autonomous
system.
